\documentclass[10pt,a4paper]{amsart}
\usepackage[margin=19mm]{geometry}
\usepackage{amsmath,amssymb,mathtools}
\usepackage[scr=rsfs,cal=euler]{mathalfa}
\usepackage{xfrac}
\usepackage[unicode,hidelinks,bookmarks=false]{hyperref}
\newcommand{\ie}{i.e.\ }
\newcommand{\cf}{cf.\ }
\newcommand{\resp}{respectively\ }
\newcommand{\Subsection}{\subsection}
\providecommand{\nobreakdash}{\nobreak\hbox{-}}
\setlength{\emergencystretch}{2em}
\multlinegap0pt
\usepackage{mdframed}
\usepackage{mdframed}
\usepackage{mdframed}
\usepackage{mdframed}
\usepackage[normalem]{ulem}
\usepackage{dsfont}
\usepackage{mathtools}
\usepackage{amssymb}
\usepackage{mdframed}
\usepackage{bm}
\usepackage{tikz}
\newtheorem{lemma}{Lemma}
\DeclareMathOperator{\Gr}{Gr}
\def \R{{\mathbb R}}
\def\S{ {\mathbb{S}} }
\def\SL{   {\rm SL}}
\newcommand{\dd}{\,\mathrm{d}}
\title{Rigidity and equidistribution of random walks by diffeomorphisms near the conservative regime}
\author{Timoth\'ee B\'enard, Zhiyuan Zhang}
\date{Source: arXiv:2605.27008v1 (CC BY 4.0)}
\begin{document}
\maketitle

\noindent\emph{Source and licence.} Rigidity and equidistribution of random walks by diffeomorphisms near the conservative regime. Version arXiv:2605.27008v1 (CC BY 4.0), distributed by arXiv under the Creative Commons Attribution 4.0 International licence, http://creativecommons.org/licenses/by/4.0/, as recorded in the arXiv licence metadata for this exact version and shown on the abstract page https://arxiv.org/abs/2605.27008v1. Full licence notice: \texttt{licenses/arxiv-2605.27008-CC-BY-4.0.txt}.\par\smallskip

\noindent\textit{Editorial scope.} Counted unit: the article's lemma mu0neighborhoodfine (source0.tex lines 3277--3279). The article's complete original proof of it is delivered with it (source0.tex lines 3291--3315, one \texttt{proof} environment of 25 lines, which the article places away from the statement). No mathematics is written, repaired or re-derived: the statement and the proof are the article's own lines, in the article's own notation, and no retained line is substituted or re-flowed.  Retained uncounted as the source-local context the proof needs: lines 3283--3285; lines 3287--3289.  Nothing was omitted from any retained span, so the package carries no omission record.  No retained line contains a citation, so the wrapper carries no bibliography and no bibliography entry.  No retained line contains a figure environment, an image-inclusion command or drawing code, so no image is delivered and the package declares no asset dependency.  Editorial formatting only, in addition: an AMS \texttt{amsart} preamble that restates the article's own declarations and macros used by the retained lines, namely the environments \texttt{lemma} on separate counters, together with the standard packages those lines need.  The retained environments print in this document as Lemma 1 in order of appearance, whereas the article prints the same items with the numbering of its own full document.  The complete native source of the article as distributed in the arXiv e-print for this exact version is bundled unchanged as \texttt{sources/201463/source0.tex}.
			\begin{align}\label{eq-cv-L1-LLN}
				n^{-1}\int_{\SL_{d}(\R)} \log \|gw\| \,\,\dd\mu_{0}^{*n}(g) \rightarrow \sigma_{d-k+1}+\dots +\sigma_{d}.
			\end{align}

			\begin{align}\label{eq-cv-norm-LLN}
				n^{-1}\int_{\SL_{d}(\R)} \log \|g\|_{\wedge^{k}\R^d} \,\,\dd\mu_{0}^{*n}(g) \rightarrow \sigma_{d-k+1}+\dots +\sigma_{d}.
			\end{align}		

		\begin{lemma}  \label{lem mu0neighborhoodfine}
			The measure $\mu_{0}$ has a $b$-gap for every $b\in \{1, \dots, d-1\}$.
		\end{lemma}

		\begin{proof}[Proof of Lemma \ref{lem mu0neighborhoodfine}] 	
				Let $b\in \{1, \dots, d-1\}$. We check that $\mu_{0}$ has a $b$-gap. 		Given a subspace $S\in \Gr(d,k)$, we write $e_{S}	\in \wedge^{k}\R^d$ a unit vector spanning the line $\wedge^{k}S$.	Let $V\in \Gr(d, d-b)$, let $u\in V^\perp$ with norm $1$.  Observe that
			\begin{align*}
				\|g e_{(V+\R u)}\|= \|g (e_{V}\wedge u)\| =\|g e_{V}\| \|e_{gV}\wedge gu\|= \|g e_{V}\|\| P_{gV^{\perp}} gu  \|.
			\end{align*}
			It follows that					
			\begin{align}\label{eq-proj-rationorm}
				\| P_{gV^{\perp}} gu  \|\leq \frac{\|g\|_{\wedge^{d-b+1}\R^d}}{\|ge_{V}\|} 
			\end{align}
			Applying \eqref{eq-cv-L1-LLN}, \eqref{eq-cv-norm-LLN}, we deduce that uniformly in $V$, 
			\begin{align}\label{eqsupE}
				\limsup_{n} n^{-1}	\int_{\SL_{d}(\R)}  \log \sup_{u \in \S(V^{\perp})} \| P_{gV^{\perp}} gu  \|  \dd\mu_{0}^{* n_0}(g) 	 \leq  \sigma_{b}		
			\end{align}
			On the other hand, by using Cartan's decomposition, one sees that
			\begin{align}\label{eq-inf-rationorm}
				\inf_{u \in \S(V)} \| gu \| \geq \frac{\|g e_{V}\|}{\|g\|_{\wedge^{d-b-1}\R^d }}.		 
			\end{align}
			Applying \eqref{eq-cv-L1-LLN}, \eqref{eq-cv-norm-LLN}, we deduce that uniformly in $V$, 
			\begin{align}\label{eqinfE}
				\liminf_{n} n^{-1} \int_{\SL_{d}(\R)} \log \inf_{u \in \S(V)} \| g v \|   \dd\mu_{0}^{* n_0}(g) 		\geq \sigma_{b+1}		
			\end{align}
			Recalling $\sigma_{b}<\sigma_{b+1}$, the combination of \eqref{eqsupE} and \eqref{eqinfE} shows that $\mu_{0}$ has a $b$-gap.
			
			
		\end{proof}

\end{document}
